% =============================================================================
% bab2-tinjauan-pustaka.tex — BAB II TINJAUAN PUSTAKA
% Skripsi / Tesis / Disertasi (Lampiran I)
% =============================================================================

\chapter{TINJAUAN PUSTAKA}
\label{bab:tinjauan-pustaka}

\section{Kajian Pustaka}
\label{sec:kajian-pustaka}

Tuliskan kajian pustaka (literature review) yang relevan dengan topik penelitian ini.
Uraikan teori-teori, konsep, dan temuan penelitian terdahulu yang menjadi landasan
teoretis penelitian ini \cite{contohSumber2025}.

\subsection{Teori Pertama}
\label{subsec:teori-pertama}

Uraikan teori pertama yang relevan. Sub-bab menggunakan format penomoran 2.1.1 dan
seterusnya. Sub-bab dapat dilanjutkan hingga 3 tingkat (contoh: 2.1.1.1).

\subsubsection{Sub-subbab Ketiga}
\label{subsubsec:sub-subbab-ketiga}

Ini adalah level ketiga dari penomoran sub-bab (2.1.1.1). Sesuai aturan FEB,
penomoran maksimal 3 tingkat.

\subsection{Teori Kedua}
\label{subsec:teori-kedua}

Uraikan teori kedua yang relevan.

\section{Penelitian Terdahulu}
\label{sec:penelitian-terdahulu}

Sajikan ringkasan penelitian terdahulu yang relevan. Dapat disajikan dalam bentuk
tabel seperti berikut:

\begin{table}[H]
\caption{Ringkasan Penelitian Terdahulu}
\label{tab:penelitian-terdahulu}
\begin{isitable}
\begin{tabular}{p{0.5cm} p{3cm} p{4cm} p{4cm} p{2.5cm}}
\toprule
\textbf{No} & \textbf{Penulis (Tahun)} & \textbf{Judul} & \textbf{Metode} & \textbf{Hasil} \\
\midrule
1 & Penulis (2023) & Judul penelitian pertama & Metode yang digunakan & Hasil utama \\
2 & Penulis (2022) & Judul penelitian kedua & Metode yang digunakan & Hasil utama \\
3 & Penulis (2021) & Judul penelitian ketiga & Metode yang digunakan & Hasil utama \\
\bottomrule
\end{tabular}
\end{isitable}
\end{table}

\section{Kerangka Pemikiran}
\label{sec:kerangka-pemikiran}

Uraikan kerangka pemikiran penelitian ini. Gambarkan hubungan antar variabel atau
konsep yang diteliti. Kerangka pemikiran dapat disajikan dalam bentuk bagan/diagram.

\begin{figure}[H]
\centering
% Ganti dengan gambar kerangka pemikiran Anda
% \includegraphics[width=0.8\textwidth]{figures/kerangka-pemikiran}
\fbox{\parbox{10cm}{\centering\vspace{2cm}[Gambar Kerangka Pemikiran]\vspace{2cm}}}
\caption{Kerangka Pemikiran Penelitian}
\label{fig:kerangka-pemikiran}
\end{figure}

\section{Hipotesis Penelitian}
\label{sec:hipotesis}

Berdasarkan kajian pustaka dan kerangka pemikiran yang telah diuraikan, hipotesis
dalam penelitian ini adalah:

\begin{enumerate}
  \item[\textbf{H\textsubscript{1}}:] Terdapat pengaruh signifikan antara variabel X terhadap variabel Y.
  \item[\textbf{H\textsubscript{2}}:] Terdapat pengaruh signifikan antara variabel Z terhadap variabel Y.
  \item[\textbf{H\textsubscript{3}}:] Variabel X dan Z secara simultan berpengaruh terhadap variabel Y.
\end{enumerate}
